\documentclass[11pt]{article}
\usepackage{amsmath,amsthm,amssymb}
\usepackage[margin=1in]{geometry}
\usepackage{booktabs}
\usepackage{array}
\usepackage{stmaryrd}
\usepackage[colorlinks=true,linkcolor=blue,citecolor=blue,urlcolor=blue]{hyperref}

\newtheorem{theorem}{Theorem}
\newtheorem{proposition}[theorem]{Proposition}
\newtheorem{corollary}[theorem]{Corollary}
\newtheorem{lemma}[theorem]{Lemma}
\theoremstyle{definition}
\newtheorem{definition}[theorem]{Definition}
\newtheorem{remark}[theorem]{Remark}
\newtheorem{example}[theorem]{Example}

\newcommand{\Cat}{\mathsf{Cat}}
\newcommand{\Set}{\mathsf{Set}}
\newcommand{\Ab}{\mathsf{Ab}}
\newcommand{\Sc}{\mathcal{S}}
\newcommand{\C}{\mathcal{C}}
\newcommand{\D}{\mathcal{D}}
\newcommand{\Oc}{\mathcal{O}}
\newcommand{\Sk}{\mathrm{Sk}}
\newcommand{\Ob}{\mathrm{Ob}}
\newcommand{\Mor}{\mathrm{Mor}}
\newcommand{\Hom}{\mathrm{Hom}}
\newcommand{\dom}{\mathrm{dom}}
\newcommand{\cod}{\mathrm{cod}}
\newcommand{\id}{\mathrm{id}}
\newcommand{\Z}{\mathbb{Z}}
\newcommand{\Aut}{\mathrm{Aut}}
\newcommand{\Inn}{\mathrm{Inn}}
\newcommand{\Out}{\mathrm{Out}}
\newcommand{\Zc}{Z}
\newcommand{\gerbe}{\mathfrak{g}}
\newcommand{\Dfr}{\mathcal{D}}
\newcommand{\bowtiemap}{\bowtie}
\newcommand{\GL}{\mathrm{GL}}

\title{A classical $H^2$ obstruction, read as supply-chain inventory:\\
\large the inventory gerbe and its band}
\author{MacBeth\thanks{Kodamai. Source, verification scripts, and this note live at
\texttt{github.com/scotmacbeth/work-in-progress}. Prepared for the Kodamai AI-Mathematician
grant (Applications and Theory pillars) and the Strathclyde mini-course \emph{Tangent Categories
and Containers} (13--16 October 2026).}\\
\small Kodamai}
\date{10 October 2026 (corrected revision of the note dated 9 October 2026)\\[2pt]
\small Circulated to Rick (reviewer) and the Kodamai group; CC Robin Langer.\\[2pt]
\small supersedes content commit \texttt{9eb0941}; content commit \texttt{PENDING}}

\begin{document}
\maketitle

\begin{abstract}
A supply chain, abstracted to a small category $\C$, carries inventory as local stock-data
glued along logistics operations. When charts are compared by \emph{equality} of counts,
consistency is a degree-one class $[\theta_R]\in H^1(N(\C);R)$ (companion note). Here we
refine the comparison to \emph{isomorphism of stock-states}: lots are fungible, so charts match
only up to a relabelling group $K=\Aut(F)$, and the gluing data becomes an $\Aut F$-gerbe over
the nerve $N(\C)$. The resulting obstruction is \emph{not new mathematics}: it is the classical
central-extension lifting class on $N(\C)$, read for the first time as an inventory invariant.
We make this precise and honest. The obstruction is \emph{staged} after Giraud: a global
inventory exists iff \emph{two} obstructions vanish --- first the \emph{band}
$\beta\in H^1(N(\C);\Out(\Aut F))$, a genuinely nonabelian datum, and then the \emph{gerbe
class} $[\gerbe]\in H^2(N(\C);Z(\Aut F))$ (Corollary~\ref{cor:aprime}). The gerbe class vanishes
on every $1$-dimensional (free / DAG / tree) chain, and is nonzero already on a single reversible
loop $B(\Z/2)$ --- but only for a \emph{nonabelian} fibre with nontrivial centre; the minimal
witness is a four-lot ring, $\Aut F=D_4$ (Theorem~\ref{thm:B}). It is logically independent of
the Zappa--Sz\'ep weld $[\omega]$, not because the two live over different coefficients --- on
$B(\Z/2)$ they can lie in the \emph{same} group and even coincide --- but because $[\gerbe]$
depends on the inventory datum and $[\omega]$ does not (Theorem~\ref{thm:C}). Finally, both
strata are invisible to plain fungibility: $\Aut F=S_n$ is centreless \emph{and} complete for
$n\ge3,\ n\ne6$, so the two-tier phenomenon is a feature of \emph{gauge/torsor-structured} stock,
not of interchangeable commodities (Theorem~\ref{thm:D}). The result places inventory consistency
as a fourth distinct $H^2$ home on $\mathrm{DCont}\simeq\Cat$, distinguished by its \emph{input
datum}.
\end{abstract}

\section{Introduction}\label{sec:intro}

Composition is where correctness is won or lost, and inventory is where the price is paid. When
two delivery routes reconverge on the same stock-keeping unit, when stock is moved and merged
and counted again downstream, the question is always whether the local records glue to one
consistent global record. When they do not, the failure has an economic name: phantom
inventory, a reconciliation that never closes, a write-off. A compositional theory of supply
chains earns its keep by telling you \emph{which} such failures can occur and \emph{where} they
live.

The categorical model is standard. A supply chain --- nodes as objects, admissible logistics
operations as morphisms --- is a small category $\C$; equivalently, by Ahman and
Uustalu~\cite{au-dccat}, a \emph{directed container}. Its nerve $N(\C)$ is the simplicial set
whose $n$-simplices are composable strings of operations, and the cohomology of $N(\C)$ is the
natural home for obstructions to gluing local data along $\C$.

In a companion note~\cite{h1note} we asked when inventory, recorded as counts transported along
operations, glues consistently. Writing the inventory as a resource presheaf $R$ on $\C$, the
answer was clean: consistency is the vanishing of a \emph{degree-one} class
$[\theta_R]\in H^1(N(\C);R)$, the route-reconvergence discrepancy of stock counts, and this
class is logically independent of the Zappa--Sz\'ep weld obstruction $[\omega]\in H^2(\Sk_\C;
\Dfr)$ that governs whether $\C$ factors as a product of two chains at all~\cite{gobstr,orch}.
That note flagged one refinement as the natural next question, and the present note answers it.

\paragraph{The gap.}
Equality of counts is the wrong comparison when lots are \emph{fungible}. Two pallets of the
same grade are interchangeable; a local stock-state is a labelled configuration, but the labels
carry no information, so two charts should be regarded as matching when they are
\emph{isomorphic}, not when they are literally equal. Each inventory fibre $F$ then carries a
\emph{relabelling group} $K:=\Aut(F)$, and comparison of local charts is canonical only up to
$K$. Gluing up-to-isomorphism is not a gluing-of-sets problem; it is a gluing-of-\emph{stacks}
problem --- a gerbe over $N(\C)$. The companion note anticipated that the obstruction would rise
to $H^2$, but left three questions open. Does the obstruction equal $[\omega]$, which also lives
in $H^2$? What exactly must vanish for a global inventory to exist? For which fibres is the
phenomenon even nonvacuous?

\paragraph{What is new, and what is not.}
We settle all three. The honest framing matters, and we state it at the outset: the obstruction
is \emph{classical mathematics}. It is the central-extension lifting class of Schreier and
Eilenberg--MacLane, specialised to the nerve $N(\C)$ and pulled back along the relabelling
transport; for small categories this is the extension theory of Ho~\cite{ho} and
Wells~\cite{wells}, classified cohomologically by Baues--Wirsching~\cite{bw}. An arXiv search
(October 2026) found no prior reading of it as an inventory invariant, so the \emph{application}
is new; the \emph{mathematics} is not. We pitch it accordingly --- a classical obstruction,
newly read as inventory --- and the work of the note is to get the reading exactly right.

\paragraph{The results.}
Fix a fibre $F$ with relabelling group $K=\Aut(F)$, centre $\Zc:=Z(K)$, inner group
$\Inn K=K/\Zc$, and outer group $\Out K$. Two sequences organise everything:
\begin{align}
&1 \longrightarrow \Zc \longrightarrow K \xrightarrow{\ \pi\ } \Inn K \longrightarrow 1
&&\text{(central extension of $K$)}\label{eq:star}\tag{$\star$}\\
&1\to\Zc\to K\xrightarrow{\ \mathrm{conj}\ }\Aut K\xrightarrow{\ q\ }\Out K\to1,\quad
\Inn K=\ker q
&&\text{(inner/outer tower)}\label{eq:tower}\tag{$\star\star$}
\end{align}
The gluing datum of fungible inventory is an $\Aut F$-gerbe over $N(\C)$ whose structure group is
$K$. Classifying it is \emph{staged}, exactly as in Giraud's non-abelian cohomology: the
\emph{band} $\beta\in H^1(N(\C);\Out K)$ --- the outer part of the relabelling transport --- must
vanish before the degree-two class is even defined, and then that class
$[\gerbe]\in H^2(N(\C);\Zc)$ is the lifting obstruction of~\eqref{eq:star}. Our results are:

\begin{theorem}[Rectification $\Leftrightarrow$ vanishing; Theorem~\ref{thm:A}]
Once the band is trivial, a global inventory --- a strict functor $\C\to BK$ realising the
up-to-iso gluing --- exists if and only if $[\gerbe]=0$ in $H^2(N(\C);\Zc)$. This is the Schreier /
Eilenberg--MacLane central-extension lifting criterion, specialised to $N(\C)$.
\end{theorem}

\begin{corollary}[The staged obstruction; Corollary~\ref{cor:aprime}]
A global inventory exists if and only if \emph{both} strata vanish: \textup{(i)} the band
$\beta=0\in H^1(N(\C);\Out K)$ and \textup{(ii)} the gerbe class $[\gerbe]=0\in H^2(N(\C);\Zc)$,
the latter well-defined only once \textup{(i)} holds.
\end{corollary}

\begin{theorem}[Degree-genuineness; Theorem~\ref{thm:B}]
If $|N(\C)|$ is homotopy equivalent to a \emph{$1$-dimensional} complex --- in particular if $\C$
is the free category on a quiver (a tree-like or DAG chain) --- then $H^{\ge2}(N(\C);\Zc)=0$, so
$[\gerbe]=0$ for every fibre. There is a minimal one-node witness, $\C=B(\Z/2)$ with the fibre a
four-lot ring ($K=\Aut F=D_4$), for which $[\gerbe]\neq0$.
\end{theorem}

\begin{theorem}[Independence from the weld; Theorem~\ref{thm:C}]
$[\gerbe]\in H^2(N(\C);\Zc)$ and the Zappa--Sz\'ep weld $[\omega]\in H^2(\Sk_\C;\Dfr)$ are
logically independent: over the family of pairs (inventory datum, weld datum) on a fixed $\C$,
all four vanishing patterns occur. The mechanism is \emph{not} a difference of coefficients ---
on $\C=B(\Z/2)$ the two classes can lie in the same group $H^2(\Z/2;\Z/2)$ and even coincide ---
but a difference of \emph{input data}: $[\gerbe]$ depends on the inventory fibre, $[\omega]$ does
not.
\end{theorem}

\begin{theorem}[Nonvacuousness; Theorem~\ref{thm:D}]
$[\gerbe]$ can be nonzero only if $K$ is \emph{nonabelian with nontrivial centre}
($\Inn K\ne1$ and $\Zc\ne1$); the minimal such $K$ are $D_4$ and $Q_8$. For $F$ a bare
$n$-element set of indistinguishable lots, $\Aut F=S_n$ is abelian or centreless for every $n$
and complete (so $\Out=1$) for $n\ge3,\ n\ne6$: fungible stock carries \emph{neither} obstruction.
The two-tier phenomenon is a feature of gauge/torsor-structured stock, not of interchangeable
commodities.
\end{theorem}

The main results are visible by page two. The picture is a \emph{two-tier} obstruction with a
\emph{staged} degree-two tier:

\begin{center}
\begin{tabular}{llll}
\toprule
tier & charts compared by & obstruction & coefficients, degree\\
\midrule
1 & equality of counts & $[\theta_R]$ & $R$ (resource group), $H^1$\\
2a (band) & iso of states, outer part & $\beta$ & $\Out(\Aut F)$ (nonabelian), $H^1$\\
2b (gerbe) & iso of states, central part & $[\gerbe]$ & $Z(\Aut F)$ (fibre centre), $H^2$\\
\midrule
(precond.) & does $\C$ factor at all & $[\omega]$ & $\Dfr$ (weld vertex groups), $H^2$\\
\bottomrule
\end{tabular}
\end{center}

\paragraph{The method.}
One classical fact, specialised to $N(\C)$, does the work. Comparing inventory charts across an
operation transports the relabelling group $K$ by an automorphism, functorial only up to inner
automorphism; its outer reduction is the band $\beta$ (Lemma~\ref{lem:band0}), and $\beta=0$ iff
the transport is inner-valued up to relabelling (Theorem~\ref{thm:band}). When $\beta=0$ the
transport descends to a functor $\rho:\C\to B(\Inn K)$, and the gerbe is its \emph{gerbe of lifts}
through~\eqref{eq:star}: choosing set-section lifts $T_f=s(\rho_f)$ exhibits a $\Zc$-valued
$2$-cocycle $\gamma$ whose class is $[\gerbe]$ (\S\ref{sec:gerbe}). Rectification then reads as
trivialising $\gamma$ by a central $1$-cochain --- an elementary coboundary computation
(Theorem~\ref{thm:A}). The remaining theorems are a simplicial computation on the nerve
(Theorem~\ref{thm:B}), a comparison with explicit witnesses (Theorem~\ref{thm:C}), and the
classical automorphism arithmetic of $S_n$ (Theorem~\ref{thm:D}).

\paragraph{Consequences and placement.}
Three points up front. (i) The refinement is \emph{exactly} one categorical level: equality-gluing
is $H^1$, iso-gluing is $H^2$, with a nonabelian $H^1$ band in between that must vanish first.
(ii) Because $[\gerbe]$ and $[\omega]$ can share a group, the temptation to identify ``inventory
is inconsistent up to relabelling'' with ``the chain fails to weld'' is a genuine trap ---
they can even be equal --- and only the \emph{input datum} separates them; Theorem~\ref{thm:C}
makes this precise. (iii) The nonvacuousness criterion sharpens the applications story: the
economically relevant tier-two case is not plain fungibility of many interchangeable units
(centreless, complete $S_n$) but \emph{gauge-valued} stock --- allocation defined up to a group
action --- whose automorphism group is nonabelian with a centre. Within the Kodamai program,
$[\gerbe]$ is the \emph{fourth} distinct $H^2$ home on $\mathrm{DCont}\simeq\Cat$, after the
Zappa--Sz\'ep holonomy, Ferri prunability, and the Baues--Wirsching class $H^2_{\Cat}$;
it is distinguished from them by its \emph{input datum} (the inventory fibre), not by a
coefficient-system discriminator.

\paragraph{Outline.}
Section~\ref{sec:prelim} recalls supply chains as directed containers, the two
sequences~\eqref{eq:star}--\eqref{eq:tower}, and the Giraud strata we use.
Section~\ref{sec:gerbe} defines the object as the gerbe of lifts and exhibits its cocycle.
Section~\ref{sec:A} proves the rectification criterion (Theorem~A) and identifies it as classical.
Section~\ref{sec:band} proves the band criterion and the staged corollary (the band boundary).
Section~\ref{sec:B} proves degree-genuineness (Theorem~B). Section~\ref{sec:C} proves independence
from the weld (Theorem~C). Section~\ref{sec:D} proves the nonvacuousness criterion and the
fungibility dichotomy (Theorem~D). Section~\ref{sec:conclusion} states the scope, the honest open
problem (coupled strata), and the two-tier bottom line.

\paragraph{Note on this revision.}
This note corrects the version of 9 October 2026 (commit \texttt{9eb0941}) following Rick's
referee report of 10 October. The corrections are substantive and we state them openly rather
than paper over them: the earlier degree-genuineness hypothesis (``homotopy $1$-truncated'')
was false and is replaced by ``homotopy equivalent to a $1$-complex''
(\S\ref{sec:B}); the earlier minimal witness ($K=\Z/2$) was wrong --- a bare set carries no
tier-two class --- and is replaced by the four-lot ring $K=D_4$ (\S\ref{sec:B},~\ref{sec:D}); the
earlier ``different coefficients'' explanation of independence was false and is replaced by the
input-data mechanism (\S\ref{sec:C}); and the band, previously hypothesised away, is now settled
(\S\ref{sec:band}). The gerbe cocycle is renamed $\gamma$ to free the symbol $\omega$ for the weld.

\section{Preliminaries}\label{sec:prelim}

The section records only what the proofs use. A reader fluent in directed containers and
non-abelian cohomology can skim to Section~\ref{sec:gerbe}.

\subsection{Supply chains as directed containers and their nerve}

A supply chain is modelled as a small category $\C$: objects are nodes (warehouses, plants,
stock locations), morphisms are admissible logistics operations, composition is sequencing. By
Ahman--Uustalu~\cite{au-dccat} a small category is the same thing as a directed container, so the
model lives natively in the container program.

The \emph{nerve} $N(\C)$ is the simplicial set with $N(\C)_0=\Ob\C$, $N(\C)_1=\Hom\C$, and
$N(\C)_n$ the set of composable strings $c_0\xrightarrow{f_1}c_1\to\cdots\xrightarrow{f_n}c_n$.
For an abelian coefficient group $A$ with (here) trivial $\C$-action, $H^\bullet(N(\C);A)$ is the
cohomology of the normalised cochain complex with the usual alternating-face differential. Two
facts are used below:
\begin{itemize}
\item For a one-object category $\C=BG$ (a group $G$ viewed as a category), $N(BG)$ is the
standard classifying space and $H^\bullet(N(BG);A)=H^\bullet(G;A)$ is group cohomology.
\item If $|N(\C)|$ is homotopy equivalent to a $1$-dimensional CW complex then
$H^{\ge2}(N(\C);A)=0$. This holds for the free category $F(Q)$ on a quiver $Q$, whose classifying
space is homotopy equivalent to the graph $|Q|$.
\end{itemize}
We emphasise the second fact in its correct form: it is \emph{homotopy $1$-dimensionality}, not
homotopy $1$-truncatedness, that forces $H^{\ge2}=0$. A $K(\pi,1)$ such as $B(\Z/2)$ is
$1$-truncated yet has $H^{\ge2}\ne0$; see Remark~\ref{rem:truncated}.

\subsection{Inventory fibres and the two sequences}\label{sec:prelim-ext}

An \emph{inventory fibre} $F$ is the local type of a stock-state at a node: a configuration of
lots. Fungibility means the labels carry no information, so the relevant symmetry is the
\emph{relabelling group}
\[
K:=\Aut(F),\qquad \Zc:=Z(K)\ \text{(centre)},\qquad \Inn K:=K/\Zc\ \text{(inner group)},
\]
fitting into the central extension~\eqref{eq:star}. We write $\pi:K\to\Inn K$ for the quotient
and reserve $s:\Inn K\to K$ for a chosen \emph{set-section} with $s(e)=e$ (not a homomorphism in
general). The band lives in the outer group $\Out K$ of the inner/outer
tower~\eqref{eq:tower}; $\Inn K=\ker q\cong K/\Zc$ is the image of conjugation.

\subsection{The gerbe, its band, and the Giraud strata}\label{sec:prelim-gerbe}

Comparing two inventory trivialisations of the same fibre across an operation $f$ identifies them
only up to the relabelling group $K$, i.e.\ transports $K$ by an \emph{automorphism} $\alpha_f\in
\Aut K$. Re-trivialising a chart conjugates the transport, so $\alpha$ is functorial only up to
inner automorphisms: $\alpha_{gf}=\iota\cdot\alpha_g\alpha_f$ for some $\iota\in\Inn K$. Collecting
the local trivialisations into a stack makes the gluing datum an $\Aut F$-gerbe over $N(\C)$: a
stack of groupoids locally equivalent to $B(\Aut F)$ (Giraud~\cite{giraud}, III; Breen~\cite{breen}
for the bitorsor/cocycle presentation). A \emph{global inventory} is a global section realised on
the nose --- a strict functor $\C\to BK$ whose transition data glue strictly.

Giraud's classification of a $G$-gerbe over a base $X$ (for us $G=\Aut F=K$, $X=N(\C)$) factors
into \emph{three strata}:

\begin{center}
\begin{tabular}{clll}
\toprule
stratum & datum & invariant & home\\
\midrule
1 (unbanded) & structure $2$-group cocycle & $\pi_0 \cong H^1(X;\mathrm{AUT}\,K)$ & nonabelian $H^1$\\
\textbf{2 (band)} & outer transport $\rho:=q\circ\alpha$ & $\beta:=[\rho]\in H^1(X;\Out K)$ & nonabelian $H^1$\\
3 (fixed band) & centre cocycle & $[\gerbe]\in H^2(X;\Zc)_\beta$ & abelian $H^2$\\
\bottomrule
\end{tabular}
\end{center}

We work the classification at strata $2$ and $3$, which is where a \emph{strict global section}
lives: stratum $1$ concerns which $2$-group presents the gerbe, which for a fixed fibre $F$ and
constant relabelling group is already pinned. The \textbf{band} $\beta$ must vanish before the
degree-two coefficients $\Zc$ are untwisted and $[\gerbe]$ is defined; this staging is the content
of Section~\ref{sec:band}. For stratum $3$ we need only one classical fact, imported rather than
re-derived: an abelian-band gerbe over $N(\C)$ with band $A$ (trivial $\C$-action) is classified
by $H^2(N(\C);A)$, the trivial class being the one admitting a global section (Giraud~\cite{giraud},
III). Everything else --- the cocycle, its class, the rectification algebra, the band criterion ---
we compute in place.

\paragraph{Conventions for nonabelian $H^1$.}
For a group $Q$ (constant coefficients) a \emph{cocycle} is a functor $\C\to BQ$, i.e.\ a map
$\gamma:\Mor\C\to Q$ with $\gamma(gf)=\gamma(g)\gamma(f)$, $\gamma(\id)=e$; two cocycles are
\emph{cohomologous} iff related by an object-relabelling $\lambda:\Ob\C\to Q$,
$\gamma'(f:x\to y)=\lambda(y)\gamma(f)\lambda(x)^{-1}$. Then $H^1(N(\C);Q):=Z^1/{\sim}$, pointed
by $\gamma\equiv e$; $[\gamma]=0$ means $\gamma\sim e$. For $Q$ abelian this is ordinary
simplicial $H^1$; for $Q$ nonabelian it is a pointed set.

\section{The inventory gerbe as the gerbe of lifts}\label{sec:gerbe}

We define the object directly as the \emph{gerbe of lifts} of the inner transport through the
central extension~\eqref{eq:star}. This is the honest construction: it makes transparent that the
degree-two class is extra lifting data pulled back along the transport, not a number computed from
the fibre $F$ in isolation.

\begin{definition}[The inventory gerbe]\label{def:gerbe}
Fix a fibre $F$ with $K=\Aut F$ and suppose the band is trivial ($\beta=0$;
Section~\ref{sec:band}), so the relabelling transport is inner-valued up to relabelling and
descends to a functor
\[
\rho:\C\longrightarrow B(\Inn K),\qquad \Inn K=K/\Zc.
\]
The \emph{inventory gerbe} is the $\Zc$-gerbe of lifts of $\rho$ through $\pi$
in~\eqref{eq:star}: its local objects are lifts of $\rho_f$ to $K$, and its band is the centre
$\Zc$ with trivial $\C$-action. A \emph{global inventory} is a global section --- a strict
functor $\C\to BK$ lifting $\rho$.
\end{definition}

\begin{remark}[Why the gerbe of lifts, and not a bare $K$-gerbe]\label{rem:lifts}
It is tempting to call the object a ``$K$-gerbe with trivial band'' and invoke the band formalism
to reduce its coefficients to $\Zc$. That is imprecise: a $K$-gerbe with trivial band and the
$\Zc$-gerbe of lifts of $\rho$ are \emph{different} objects --- for $\rho=\id$ the inner form
$\mathbf{Tors}(\Inn K)$ has a global section while the $\Zc$-gerbe of lifts does not --- so
``global section $\Leftrightarrow[\gerbe]=0$'' depends on which one is meant. We therefore
\emph{define} the object as the gerbe of lifts (equivalently, as the $\mathrm{AUT}(F)$-pseudofunctor
of Remark~\ref{rem:kinvariant}), and no separate band-is-the-centre lemma is needed: the
coefficients are $\Zc$ by construction.
\end{remark}

\subsection{The cocycle}

\begin{definition}[Section cocycle]\label{def:cocycle}
Choose set-section lifts $T_f:=s(\rho_f)\in K$ for each operation $f$, with $T_{\id}=e$. Because
both $T_gT_f$ and $T_{gf}$ map to $\rho_g\rho_f=\rho_{gf}$ under $\pi$, their ratio lands in
$\ker\pi=\Zc$:
\begin{equation}\label{eq:defect}
T_gT_f=\gamma(g,f)\cdot T_{gf},\qquad \gamma(g,f)\in\Zc.
\tag{$\dagger$}
\end{equation}
Thus $\gamma=\rho^*(\text{extension $2$-cocycle of~\eqref{eq:star}})$ is a normalised
$2$-cochain on $N(\C)$ valued in $\Zc$.
\end{definition}

\begin{lemma}[{$\gamma$ is a cocycle and $[\gerbe]$ is well-defined}]\label{lem:cocycle}
Associativity of multiplication in $K$, together with the centrality of $\Zc$ (so conjugation acts
trivially on the values of $\gamma$), gives the normalised $2$-cocycle identity
\[
\gamma(h,g)\,\gamma(hg,f)=\gamma(g,f)\,\gamma(h,gf),\qquad\text{i.e.\ } \delta\gamma=0,
\]
and changing the set-section $s$ replaces $\gamma$ by a cohomologous cochain. Hence
$[\gerbe]:=\rho^*[\star]=[\gamma]\in H^2(N(\C);\Zc)$ is well-defined.
\end{lemma}

\begin{proof}
Expand $T_h(T_gT_f)=(T_hT_g)T_f$ using~\eqref{eq:defect} three times and cancel the common
$T_{hgf}$; centrality lets the $\Zc$-values commute past the $T$'s. Replacing $s$ by $s'$ sets
$T'_f=c_fT_f$ with $c_f\in\Zc$, and~\eqref{eq:defect} gives
$\gamma'(g,f)=\gamma(g,f)\,(\delta c)(g,f)$.
\end{proof}

\begin{remark}[{Face 1 is definitional: $[\gerbe]$ is a $k$-invariant, not a number from $F$}]%
\label{rem:kinvariant}
When $K$ is abelian, $\Inn K=1$, $\rho$ is trivial, and the gerbe of lifts of
Definition~\ref{def:gerbe} is \emph{trivial}: $[\gerbe]=0$. So an abelian fibre carries no
inventory class by this construction. The abelian $K$-gerbes presented by a free $2$-cocycle
$\gamma\in Z^2(N(\C);K)$ --- the ``abelian face'' of the earlier version --- are therefore
\emph{not} determined by $F$: a pseudofunctor into the $1$-category $BK$ is just a functor and
carries no $2$-cell data, so a nonzero class is extra input. Precisely, the datum is a
pseudofunctor into the automorphism $2$-group $\mathrm{AUT}(F)$ (with $\pi_1=\Out K$, $\pi_2=\Zc$),
and $[\gerbe]$ is its \emph{$k$-invariant}~\cite{bl,js}. We record this honestly: Face~1 is the
\emph{definition} of a classification datum, not a computation from the fibre. The constructive,
inventory-driven class is the gerbe of lifts of Definition~\ref{def:gerbe}, which requires a
\emph{nonabelian} $K$ (so that $\rho$ can be nontrivial) --- as Theorem~\ref{thm:D} makes precise.
\end{remark}

\section{Rectification is the vanishing of $[\gerbe]$}\label{sec:A}

\begin{theorem}[Theorem~A]\label{thm:A}
With the band trivial, a global inventory --- a strict functor $\C\to BK$ lifting $\rho$ --- exists
if and only if $[\gerbe]=0$ in $H^2(N(\C);\Zc)$, where $\Zc=Z(\Aut F)$.
\end{theorem}

\begin{proof}
Re-choosing the lift at each operation $f$ multiplies $T_f$ by an element $c_f\in\Zc$ (it moves
$T_f$ within its $\Zc$-torsor of lifts of $\rho_f$). A candidate global inventory is a $1$-cochain
$c:\Mor\C\to\Zc$ producing new transitions $T'_f=c_f T_f$, and strictness is $T'_{gf}=T'_gT'_f$.
Using~\eqref{eq:defect} and centrality,
\[
c_{gf}\,T_{gf}=T'_{gf}\overset{!}{=}T'_gT'_f=c_gc_f\,T_gT_f=c_gc_f\,\gamma(g,f)\,T_{gf},
\]
which holds iff $\gamma(g,f)=c_{gf}\,c_g^{-1}c_f^{-1}=(\delta c^{-1})(g,f)$. Hence a global
inventory exists iff $\gamma$ is a $\Zc$-coboundary, i.e.\ iff $[\gerbe]=[\gamma]=0$.
\end{proof}

\begin{remark}[This is the classical lifting criterion]\label{rem:classical}
Theorem~\ref{thm:A} is \emph{not} new. It is the Schreier / Eilenberg--MacLane criterion for
lifting a map into a quotient $\Inn K$ through a central extension~\eqref{eq:star}, here pulled back
along $\rho$ and specialised to the simplicial site $N(\C)$: for small categories this is the
extension theory of Ho~\cite{ho} and Wells~\cite{wells}, with the cohomological classification of
Baues--Wirsching~\cite{bw}. We present it as a specialisation so read, not as a theorem of this
note; the contribution is the \emph{reading} of $\rho$, $\Zc$, and $[\gerbe]$ as inventory transport,
fibre centre, and rectification obstruction.
\end{remark}

\begin{remark}[Concrete check: $D_4$ and $Q_8$]\label{rem:d4q8}
Take $\C=B((\Z/2)^2)$ and $K\in\{D_4,Q_8\}$, each with $\Zc=\Z/2$, $\Inn K=(\Z/2)^2$, and
$\rho=\id$ (the inner transport is the full inner group). A strict lift would be a splitting
homomorphism $(\Z/2)^2\to K$ of~\eqref{eq:star}; none exists, since $K$ is nonabelian while
$(\Z/2)^2$ is abelian. Correspondingly the section-cocycle $\gamma$ of~\eqref{eq:defect} is not a
coboundary, so $[\gerbe]\neq0$ and rectification fails. The non-splitting and $[\gerbe]\neq0$ are
confirmed by brute-force enumeration over $\mathbb F_2$ (\texttt{inventory\_gerbe\_h2.py}; the
reduced-nerve $H^2$ engine of the Lean witness file~\cite{leanwit}).
\end{remark}

\section{The band, and the staged obstruction}\label{sec:band}

The earlier version hypothesised the band away. We now settle it, following Giraud's three-stage
factorisation over the simplicial site $N(\C)$. This is the program's first genuinely
\emph{nonabelian} content: the band lives in $H^1$ with the nonabelian coefficient \emph{group}
$\Out K$.

\begin{definition}[Band]
The \emph{band} (lien) of the inventory gerbe is the outer reduction of the relabelling transport,
\[
\rho_{\mathrm{out}}:=q\circ\alpha:\Mor\C\longrightarrow\Out K,
\]
where $\alpha:\Mor\C\to\Aut K$ is the transport (functorial up to inner) and $q$ is as
in~\eqref{eq:tower}. Its class is $\beta:=[\rho_{\mathrm{out}}]\in H^1(N(\C);\Out K)$.
\end{definition}

\begin{lemma}[The band is always a strict cocycle]\label{lem:band0}
$\rho_{\mathrm{out}}$ is a genuine functor $\C\to B(\Out K)$, so $\beta$ is well-defined
regardless of the up-to-inner slack in $\alpha$.
\end{lemma}

\begin{proof}
Write $\alpha_{gf}=\iota\cdot\alpha_g\alpha_f$ with $\iota\in\Inn K=\ker q$. Apply $q$:
$\rho_{\mathrm{out}}(gf)=q(\iota)\,q(\alpha_g)q(\alpha_f)=\rho_{\mathrm{out}}(g)\rho_{\mathrm{out}}(f)$,
since $q(\iota)=e$. Normalisation is immediate; the inner slack dies in $\Out K$.
\end{proof}

\begin{theorem}[Band-triviality criterion; Theorem~(a)]\label{thm:band}
$\beta=0\in H^1(N(\C);\Out K)$ if and only if the transport $\alpha$ is, up to an
$\Out K$-relabelling of objects, inner-valued: iff there is a cocycle
$\tilde\alpha:\C\to B(\Aut K)$, cohomologous to $\alpha$ via an $\Aut K$-valued object relabelling,
with $\tilde\alpha(f)\in\Inn K$ for every $f$. In particular, if $\alpha$ already lands in $\Inn K$
then $\beta=0$; this is the literal ``trivial band'', and the theorem says $\beta=0$ is its
gauge-invariant closure.
\end{theorem}

\begin{proof}
($\Leftarrow$) Let $\tilde\alpha=\mu\cdot\alpha$ with $\mu:\Ob\C\to\Aut K$ and
$\tilde\alpha(f)\in\Inn K=\ker q$ for all $f$. Set $\lambda:=q\circ\mu:\Ob\C\to\Out K$. For
$f:x\to y$,
\[
\rho_{\mathrm{out}}(f)=q(\alpha(f))=q\big(\mu(y)^{-1}\tilde\alpha(f)\mu(x)\big)
=\lambda(y)^{-1}\,q(\tilde\alpha(f))\,\lambda(x)=\lambda(y)^{-1}\lambda(x),
\]
so $\lambda$ trivialises $\rho_{\mathrm{out}}$ and $\beta=0$.

($\Rightarrow$) If $\beta=0$ there is $\lambda:\Ob\C\to\Out K$ with
$\rho_{\mathrm{out}}(f)=\lambda(y)^{-1}\lambda(x)$. As $q$ is surjective, choose lifts
$\mu(x)\in\Aut K$ with $q(\mu(x))=\lambda(x)$, and set $\tilde\alpha(f):=\mu(y)\alpha(f)\mu(x)^{-1}$.
Then $q(\tilde\alpha(f))=\lambda(y)\rho_{\mathrm{out}}(f)\lambda(x)^{-1}=e$, so
$\tilde\alpha(f)\in\Inn K$; and $\tilde\alpha$ is a cocycle cohomologous to $\alpha$ via $\mu$.
\end{proof}

This is Giraud stratum $2$ specialised to $N(\C)$: the band vanishes exactly when the transport is
inner-valued up to coboundary. The proof is elementary and Schreier-style, using only the
tower~\eqref{eq:tower} and the nerve's $H^1$. For $K=D_4$ ($\Aut=D_4$ of order $8$,
$\Inn=(\Z/2)^2$, $\Out=\Z/2$ --- both parts nontrivial) over $\C=B(\Z)$ the criterion holds for all
$8$ automorphisms (\texttt{scratch/clause\_a\_D4\_check.py}).

\begin{corollary}[The full obstruction is staged; Corollary~(a$'$)]\label{cor:aprime}
A global inventory exists if and only if \emph{both} \textup{(i)} $\beta=0$ (band trivial,
Theorem~\ref{thm:band}) and \textup{(ii)} $[\gerbe]=0\in H^2(N(\C);\Zc)$ (the gerbe class, defined
once (i) holds). The obstruction is genuinely staged: $[\gerbe]$ is defined only after $\beta=0$,
otherwise the stratum-$3$ coefficients are $\beta$-twisted. Thus Theorem~\ref{thm:A} is exactly the
band-trivial stratum, and Theorem~\ref{thm:band} supplies the prior stratum.
\end{corollary}

\begin{proof}
A global inventory is a relabelling $\mu:\Ob\C\to\Aut K$ trivialising $\alpha$ outright; projecting
by $q$ it trivialises $\rho_{\mathrm{out}}$, forcing $\beta=0$, so (i) is necessary. Given (i),
Theorem~\ref{thm:band} makes $\alpha$ inner-valued, $\alpha=\mathrm{conj}\circ\rho$ for
$\rho:\C\to B(\Inn K)$; the gerbe is then the gerbe of lifts of $\rho$, neutral iff $[\gerbe]=0$ by
Theorem~\ref{thm:A}. So (i)$\wedge$(ii) is necessary and sufficient.
\end{proof}

\begin{remark}[The band is where the nonabelian content lives]\label{rem:band-nonabelian}
The band is the first obstruction in the program with a \emph{nonabelian coefficient group}
$\Out(\Aut F)$. This is an honest structural statement --- genuinely different in kind from the
abelian $H^2$ classes --- and it is what the earlier ``distinguished by its coefficient system''
slogan should have said (that slogan was false for $[\gerbe]$ versus $[\omega]$;
Theorem~\ref{thm:C}). A nontrivial-band witness, with a clean separation of the two strata:
\end{remark}

\begin{proposition}[A nonvacuous band; Proposition~(c)]\label{prop:bandwitness}
Let $K=(\Z/p)^2$ (abelian gauge group of stock) and $\C=B(\Z)$ (one reversible/cyclic operation).
For abelian $K$, $\Inn K=1$ and $\Out K=\Aut K=\GL_2(\mathbb F_p)$, so the band sees the entire
transport: $\rho_{\mathrm{out}}=\alpha$. Set $\alpha(t):=A$ for some $A\ne I$. Then
\[
H^1(N(B(\Z));\Out K)=\{\text{conjugacy classes of }\GL_2(\mathbb F_p)\},\qquad \beta=[A]\neq0.
\]
Moreover the strata separate cleanly: since $N(B(\Z))\simeq S^1$ is $1$-dimensional,
$H^2(N(\C);\Zc)=0$, so the gerbe class is \emph{forced} to zero while the band is the sole, nonzero
obstruction --- living in nonabelian $H^1$ with coefficients $\GL_2(\mathbb F_p)$, prior to and
independent of any $H^2$ story.
\end{proposition}

\begin{proof}
$N(B(\Z))\simeq S^1$; a cocycle is determined by $\rho_{\mathrm{out}}(t)=A$ and object-relabelling
by $g\in\Out K$ sends $A\mapsto gAg^{-1}$, so $H^1=\{\text{conjugacy classes}\}$ and $\beta=0$ iff
$A=I$. For $p=2$, $\GL_2(\mathbb F_2)\cong S_3$; the transvection
$A=\left(\begin{smallmatrix}1&1\\0&1\end{smallmatrix}\right)\ne I$ gives $\beta\ne0$. Clean
separation: $H^2(S^1;-)=0$ for any coefficients
(\texttt{scratch/band\_boundary\_verify.py}).
\end{proof}

\paragraph{Interpretation.}
Stock is defined up to a $(\Z/p)^2$-gauge (say two independent mod-$p$ ledger offsets); a cyclic
refurbishment process $t$ relabels the gauge by a fixed $\GL_2(\mathbb F_p)$ element $A$. If
$A\ne I$, no global choice of gauge is consistent around the loop --- a \emph{band} obstruction, an
outer twist of the structure group, strictly prior to the degree-two gerbe class. Fungibility can
never produce it (Theorem~\ref{thm:D}); gauge-structured stock does, minimally. Activating both
strata at once needs a base with $H^2\ne0$ and a nontrivial outer transport (e.g.\ $\C=B(\Z^2)$,
$N\simeq T^2$); the twisted coupling of band to centre on such a base is the natural follow-up and
is \emph{not} treated here (Section~\ref{sec:conclusion}).

\section{$[\gerbe]$ is genuinely a degree-two datum}\label{sec:B}

\begin{theorem}[Theorem~B]\label{thm:B}
\emph{(i)} If $|N(\C)|$ is homotopy equivalent to a $1$-dimensional complex --- in particular if
$\C$ is the free category on a quiver (a tree-like or DAG supply chain) --- then
$H^{\ge2}(N(\C);\Zc)=0$, so $[\gerbe]=0$ for every fibre $F$. \emph{(ii)} There is a minimal
one-node witness, $\C=B(\Z/2)$ with the fibre a four-lot ring ($K=\Aut F=D_4$), for which
$[\gerbe]\neq0$.
\end{theorem}

\begin{proof}[Proof of (i)]
If $|N(\C)|\simeq X$ with $X$ a $1$-dimensional CW complex, then $H^{\ge2}(N(\C);\Zc)=H^{\ge2}(X;\Zc)
=0$. For the free category $F(Q)$ on a quiver $Q$, $|N(F(Q))|\simeq|Q|$, a graph; hence
$H^{\ge2}=0$. It is worth seeing this directly on the reconvergent \emph{diamond}
$a\rightrightarrows\{b,c\}\rightrightarrows d$ --- the paradigmatic two-route chain --- whose nerve
is \emph{not} literally $1$-dimensional: it carries two nondegenerate $2$-simplices, the composites
$p;r$ and $q;s$ of the two routes, yet $|N|$ is still homotopy equivalent to a graph. A direct
simplicial computation over $\mathbb F_2$ (\texttt{diamond\_nerve\_H2.py}) gives chain dimensions
$(C^0,C^1,C^2)=(4,6,2)$, boundary ranks $\mathrm{rk}\,d^0=3$, $\mathrm{rk}\,d^1=2$, and hence
\[
H^0=1,\qquad H^1=1,\qquad H^2=0,\qquad \chi=4-6+2=0.
\]
The lone $H^1$ is the tier-one reconvergence loop $[\theta_R]$; $H^2=0$ says there is no gerbe on a
DAG, even one with genuine $2$-simplices.
\end{proof}

\begin{proof}[Proof of (ii)]
The fibre is four lots arranged in a ring (the $4$-cycle); its relabelling group is the dihedral
group $K=\Aut F=D_4$ of order $8$, with centre $\Zc=\{1,r^2\}\cong\Z/2$ (the half-turn) and inner
group $\Inn K=(\Z/2)^2$. Let $\C=B(\Z/2)$, with reversible operation $g$ ($g^2=\id$). Choose the
inner transport $\rho:B(\Z/2)\to B(\Inn K)$ sending $g$ to $[r]\in\Inn K$ (the class of the
rotation; $[r]$ has order $2$ since $r^2\in\Zc$). Take the set-section lift $T_g:=r$ and
$T_{\id}:=e$. Then~\eqref{eq:defect} gives
\[
\gamma(g,g)=T_gT_g\,T_{g^2}^{-1}=r^2\,T_{\id}^{-1}=r^2\neq1\in\Zc.
\]
Every normalised $1$-cochain $c:\Mor\C\to\Zc$ has $(\delta c)(g,g)=c_{g^2}c_g^{-1}c_g^{-1}=c_g^{-2}
=1$ (as $\Zc=\Z/2$ and $g^2=\id$), so $\gamma$ is \emph{not} a coboundary and $[\gerbe]$ is the
generator of $H^2(\Z/2;\Z/2)=\Z/2$. The band is trivial here ($\rho$ is inner-valued, so
$\beta=0$), and \texttt{inventory\_gerbe\_h2.py} / the Lean witness~\cite{leanwit} confirm
$[\gerbe]\neq0$.
\end{proof}

\begin{remark}[Why not $K=\Z/2$ --- homotopy $1$-truncatedness is not enough]\label{rem:truncated}
The earlier version used the \emph{abelian} fibre $K=\Z/2$ on $B(\Z/2)$ and claimed $[\gerbe]$ was
the generator of $H^2(\Z/2;\Z/2)$. This is wrong twice over. First, an abelian fibre carries
$[\gerbe]=0$ by Remark~\ref{rem:kinvariant} ($\Inn K=1$, $\rho$ trivial, no lifting gerbe): the
group $H^2(\Z/2;\Z/2)=\Z/2$ is nonzero, but the \emph{inventory class} in it is zero. Second, the
earlier degree-genuineness hypothesis ``$N(\C)$ homotopy $1$-truncated'' is itself refuted by this
very space: $B(\Z/2)=K(\Z/2,1)$ \emph{is} $1$-truncated, yet $H^2(\Z/2;\Z/2)\ne0$. The correct
hypothesis in (i) is homotopy \emph{$1$-dimensionality}. The genuine minimal witness is the
nonabelian $D_4$ of (ii).
\end{remark}

Thus $[\gerbe]$ lives at cohomological degree $2$ (nerve level $N_3$), vanishing on every
$1$-dimensional chain yet already nonzero on a single reversible loop carrying a nonabelian fibre.
For the record, $H^2((\Z/2)^2;\Z/2)=\mathbb F_2^3$: $(\Z/2)^2$ hosts three independent
$\Z/2$-gerbes --- the $D_4$ and $Q_8$ classes and the $\Z/4\times\Z/2$ class --- of which the first
were used in Remark~\ref{rem:d4q8}. (As quadratic forms on $(\Z/2)^2$ there is a single $Q_8$ class
$x^2+xy+y^2$ and three $D_4$ classes; a choice of three representatives spanning $\mathbb F_2^3$ is
$x^2,\ x^2+xy,\ x^2+xy+y^2$~\cite{leanwit}.)

\section{Independence from the Zappa--Sz\'ep weld}\label{sec:C}

\begin{theorem}[Theorem~C]\label{thm:C}
Over the family of pairs (inventory datum, weld datum) on a fixed category $\C$, the gerbe class
$[\gerbe]\in H^2(N(\C);\Zc)$ and the Zappa--Sz\'ep weld $[\omega]\in H^2(\Sk_\C;\Dfr)$ are logically
independent: all four vanishing patterns occur, including $[\gerbe]=[\omega]\neq0$. The mechanism is
a difference of \emph{input data}, not of coefficients: $[\gerbe]$ depends on the inventory fibre
$F$, which is external to $\C$, while $[\omega]$ depends only on $\C$ and a candidate factorisation.
\end{theorem}

The quantifier matters, and the earlier ``different coefficient systems'' explanation was false, so
we state the mechanism carefully before proving it. On \emph{some} categories the two classes live
in \emph{different} groups (different sites, different coefficients); on others --- notably
$\C=B(\Z/2)$ --- they live in the \emph{same} group $H^2(\Z/2;\Z/2)=\Z/2$ and can be equal. Equal
coefficients therefore cannot be what makes them independent. What makes them independent is that
they are functions of \emph{different arguments}: fixing $\C$ and varying the inventory fibre moves
$[\gerbe]$ while leaving $[\omega]$ fixed, and varying the factorisation datum moves $[\omega]$ while
leaving $[\gerbe]$ fixed. We exhibit this on a single $\C$.

\begin{proof}
Work on $\C=B(\Z/2)$, where both classes can live in $H^2(\Z/2;\Z/2)=\Z/2$. The inventory datum is
a fibre $F$; take $F\in\{D_4\text{-ring},\ \mathrm{pt}\}$, giving $[\gerbe]\neq0$
(Theorem~\ref{thm:B}(ii)) resp.\ $[\gerbe]=0$ ($\Aut F=1$). The weld datum is a candidate
Zappa--Sz\'ep factorisation of $\C$; take the re-entrant weld of~\cite{orch}, whose class
$[\omega]=\varepsilon$ is the generator of $H^2(\Z/2;\Z/2)$, resp.\ the trivial factorisation
$\Dfr=0$ with $[\omega]=0$. These two binary choices are independent --- one is a function of $F$,
the other of the factorisation --- so all four combinations are realised:

\begin{center}
\begin{tabular}{lccc}
\toprule
on $\C=B(\Z/2)$ & weld $=$ trivial ($[\omega]=0$) & weld $=$ re-entrant ($[\omega]\neq0$)\\
\midrule
fibre $=\mathrm{pt}$ ($[\gerbe]=0$) & $(0,0)$ & $(0,\neq0)$\\
fibre $=D_4$-ring ($[\gerbe]\neq0$) & $(\neq0,0)$ & $(\neq0,\neq0),\ [\gerbe]=[\omega]$\\
\bottomrule
\end{tabular}
\end{center}

In the bottom-right cell the two classes are \emph{equal} (both the generator of the same group),
yet they remain independent as functions: perturbing the fibre to $\mathrm{pt}$ kills $[\gerbe]$
without touching $[\omega]$, and vice versa. Hence neither vanishing implies the other.
\end{proof}

\begin{remark}[Not a coefficient discriminator]\label{rem:no-discriminator}
The earlier claim that ``the coefficient system is the only discriminator'' and that ``there is no
natural map comparing $H^2(N(\C);\Zc)$ with $H^2(\Sk_\C;\Dfr)$'' is withdrawn: on $B(\Z/2)$ the two
\emph{are} the same group, with the same normalised cocycles, and the comparison map is the
identity. The honest statement is the input-data one above. This is why, within the program, we
describe $[\gerbe]$ as distinguished by its \emph{input datum} (the inventory fibre) and not by a
coefficient-system discriminator.
\end{remark}

\begin{remark}[Relation to tier one needs witnesses, not a degree count]\label{rem:tier1}
It is tempting to argue that the tier-one class $[\theta_R]\in H^1$ and the tier-two class
$[\gerbe]\in H^2$ are ``immediately independent because they differ in degree''. That inference is
invalid: classes of different degree can be functionally dependent (Bocksteins, cup squares, and
$d_3$-type differentials all link classes across degrees). Independence of $[\theta_R]$ and
$[\gerbe]$ must be established by witnesses, exactly as for Theorem~\ref{thm:C}. The full inventory
obstruction is a two-tier filtration with associated graded $H^1(N(\C);R)$ (counts) and
$H^2(N(\C);\Zc)$ (relabellings); whether the tiers split or form a Postnikov-type extension (a
$d_3$ coupling $H^1\to H^3$) is not determined here (Section~\ref{sec:conclusion}).
\end{remark}

\section{When the gerbe is nonvacuous, and fungibility is doubly vacuous}\label{sec:D}

A nonzero gerbe class is a genuine --- and slightly surprising --- constraint on \emph{which}
fibres carry a tier-two obstruction, and the band adds a parallel constraint. Together they make
plain fungibility invisible to both.

\begin{theorem}[Theorem~D]\label{thm:D}
\emph{(i)} $[\gerbe]$ can be nonzero only if $K=\Aut F$ is \emph{nonabelian with nontrivial centre}:
$\Inn K\neq1$ (so $\rho$ can be nontrivial) and $\Zc\neq1$ (so the coefficients are nonzero). The
minimal such $K$ are $D_4$ and $Q_8$. \emph{(ii)} For $F$ a bare $n$-element set of indistinguishable
lots, $\Aut F=S_n$: fungible stock carries \emph{neither} stratum for all $n\ge3,\ n\ne6$, since
$S_n$ is centreless ($\Zc=1\Rightarrow[\gerbe]=0$) and complete ($\Out S_n=1\Rightarrow\beta=0$).
The low-order exceptions are disjoint: $n=6$ is the lone band exception ($\Out S_6=\Z/2$) and $n=2$
the lone centre case ($S_2=\Z/2$) --- but $S_2$ is abelian, so even there $[\gerbe]=0$ by (i).
\end{theorem}

\begin{proof}
(i) If $K$ is abelian then $\Inn K=1$, $\rho$ is trivial, and the gerbe of lifts is trivial, so
$[\gerbe]=0$ (Remark~\ref{rem:kinvariant}). If $\Zc=1$ then $H^2(N(\C);\Zc)=0$, so $[\gerbe]=0$. Thus
a nonzero class requires both $\Inn K\ne1$ and $\Zc\ne1$, i.e.\ $K$ nonabelian with nontrivial
centre; the smallest are $D_4,Q_8$ (order $8$). (ii) $S_n$ is complete for $n\ne2,6$
(centreless with all automorphisms inner, so $\Out=1$) and $\Out S_6=\Z/2$ (Rotman~\cite{rotman},
Thms.\ 7.5--7.7); $\Zc(S_n)=1$ for $n\ge3$, $=\Z/2$ for $n=2$, $=1$ for $n\le1$. Feed these into
(i) and Theorem~\ref{thm:band}. For $n=2$, $S_2=\Z/2$ is abelian, so although the \emph{group}
$H^2(\Z/2;\Z/2)$ is the home of a nonzero class, the \emph{inventory} class is $0$ by (i). The
exceptional values $n\in\{2,6\}$ are distinct and each triggers at most one stratum's coefficient
group, as tabulated (\texttt{scratch/band\_boundary\_verify.py}, \texttt{out\_s6.py}).
\end{proof}

\paragraph{The applications boundary is sharp.}
The nonabelian inventory phenomenon is \emph{invisible to fungibility}: $\ge3$ interchangeable units
create no degree-two relabelling obstruction ($S_n$ centreless) and no band obstruction ($S_n$
complete), except the isolated $S_6$ curiosity. Both nontrivial strata require
\emph{gauge/torsor-structured} stock --- a fibre whose automorphism group is nonabelian with
nontrivial centre (tier-2) or nontrivial outer group (band). This is the economically meaningful
regime: allocation defined up to a gauge group $K$ (so $\Aut F=K$ acts by translation, $\Zc=Z(K)$,
$\Out$ measuring how the process twists the gauge), not plain interchangeability. The grant headline
is exactly this: \emph{fungible inventory carries no nonabelian obstruction; the phenomenon needs
structured assets, not commodities.}

\section{Scope, open problem, and the two-tier bottom line}\label{sec:conclusion}

Refining inventory comparison from \emph{equality of counts} to \emph{isomorphism of stock-states}
raises the consistency obstruction by exactly one categorical level, through a staged obstruction.
A global inventory exists iff the band $\beta\in H^1(N(\C);\Out(\Aut F))$ vanishes \emph{and} then
the gerbe class $[\gerbe]\in H^2(N(\C);Z(\Aut F))$ vanishes (Corollary~\ref{cor:aprime}),
the latter being the classical Schreier / Eilenberg--MacLane central-extension lifting criterion
read as inventory rectification (Theorem~\ref{thm:A}, Remark~\ref{rem:classical}). The gerbe class
is genuinely degree-two, vanishing on every $1$-dimensional chain but nonzero already on a single
reversible loop $B(\Z/2)$ carrying a nonabelian four-lot ring $D_4$ (Theorem~\ref{thm:B}); it is
independent of the weld $[\omega]$ by a difference of \emph{input data}, not coefficients --- the
two can even coincide on $B(\Z/2)$ (Theorem~\ref{thm:C}); and both strata are invisible to plain
fungibility, requiring gauge/torsor-structured stock (Theorem~\ref{thm:D}). Within the container
program this is the fourth distinct $H^2$ home on $\mathrm{DCont}\simeq\Cat$, after Zappa--Sz\'ep
holonomy, Ferri prunability, and the Baues--Wirsching class --- distinguished by its \emph{input
datum}.

\paragraph{The honest open problem.}
Proposition~\ref{prop:bandwitness} witnesses a nonvacuous band but does \emph{not} classify
nontrivial-band gerbes; that is Giraud's full nonabelian machinery (the torsor over the twisted
$H^2(N(\C);\Zc)_\beta$), out of scope here. The natural next target is the \emph{coupled} stratum,
where band and centre interact: a base with $H^2\ne0$ and a nontrivial outer transport, e.g.\
$\C=B(\Z^2)$ ($N\simeq T^2$) with $K=(\Z/p)^2$ and a commuting outer pair $(A,B)$, computing the
twisted differential $d_2$ coupling $\beta$ to the centre class. We do not have this; it is
recorded as the PROVE/LEAN follow-up (a bar-complex oracle on $T^2$). Relatedly
(Remark~\ref{rem:tier1}), the coupling between tiers one and two --- split or Postnikov extension ---
is stated as a filtration but not computed.

\paragraph{For the grant.}
The worked economic example to develop is the torsor-fibre case $\Aut F=K$ acting on itself, with
$\Zc=Z(K)$ --- gauge-valued allocation, where stock is defined only up to a group action and the
tier-two obstruction has real content. The applications payoff of Theorem~\ref{thm:D} is that the
economically relevant non-rectifiable inventory is not the many-interchangeable-units case but the
gauge-structured one; and the band supplies a strictly prior, genuinely nonabelian invariant of the
same situation.

\begin{thebibliography}{99}

\bibitem{au-dccat}
D.~Ahman, T.~Uustalu,
\emph{Directed containers as categories},
Proc.\ MSFP 2016, EPTCS 207, 2016, 89--98; arXiv:1604.01187.

\bibitem{h1note}
MacBeth,
\emph{Supply-chain inventory consistency is a degree-one class, and is independent of the
Zappa--Sz\'ep weld obstruction}
(Kodamai internal note, 6 October 2026; proved).

\bibitem{gobstr}
MacBeth,
\emph{The global-closure obstruction $(G)$ is a cohomology class}
(Kodamai internal note, 11 June 2026; proved, Lean-adjacent).

\bibitem{orch}
MacBeth,
\emph{The re-entrancy obstruction is the token-mutation bit: an analytic proof for orchestration
Zappa--Sz\'ep products}
(Kodamai internal note, 20 July 2026; proved, $[\omega]=\varepsilon$ core Lean-verified).

\bibitem{twosites}
MacBeth,
\emph{Two $[\omega]$-sites are irreducibly distinct}
(Kodamai internal note, 14 August 2026; proved).

\bibitem{leanwit}
MacBeth,
\emph{Reduced-nerve $H^2$ witness engine} (Lean~4; one engine, three machine-checked witnesses:
$H^2(B(\Z/2))=\Z/2$, the $(\Z/2)^2$ diamond basis, the $D_4$ non-splitting),
Kodamai internal Lean note, 9 October 2026.

\bibitem{giraud}
J.~Giraud,
\emph{Cohomologie non ab\'elienne},
Grundlehren der mathematischen Wissenschaften \textbf{179}, Springer, 1971.

\bibitem{breen}
L.~Breen,
\emph{Bitorseurs et cohomologie non ab\'elienne},
in \emph{The Grothendieck Festschrift, Vol.\ I}, Progr.\ Math.\ \textbf{86}, Birkh\"auser, 1990,
401--476.

\bibitem{ho}
C.-W.~Ho,
\emph{A note on the classification of extensions of categories}
(central-extension lifting for small categories), 1974.

\bibitem{wells}
C.~Wells,
\emph{Extension theories for categories} (preprint, 1979); see also
D.~Bourn, \emph{The shift functor and the comprehensive factorization for internal groupoids}.

\bibitem{bw}
H.-J.~Baues, G.~Wirsching,
\emph{Cohomology of small categories},
J.\ Pure Appl.\ Algebra \textbf{38} (1985), 187--211.

\bibitem{rotman}
J.~J.~Rotman,
\emph{An Introduction to the Theory of Groups}, 4th ed.,
Graduate Texts in Mathematics \textbf{148}, Springer, 1995 (Thms.\ 7.5--7.7: $S_n$ complete for
$n\ne2,6$; $\Out S_6=\Z/2$).

\bibitem{bl}
J.~C.~Baez, A.~D.~Lauda,
\emph{Higher-dimensional algebra V: 2-groups},
Theory Appl.\ Categ.\ \textbf{12} (2004), 423--491.

\bibitem{js}
N.~Johnson, D.~Yau,
\emph{2-Dimensional Categories},
Oxford University Press, 2021.

\bibitem{weibel}
C.~A.~Weibel,
\emph{An Introduction to Homological Algebra},
Cambridge Studies in Advanced Mathematics \textbf{38}, Cambridge Univ.\ Press, 1994.

\end{thebibliography}

\end{document}
